\section{Simulation, Testing and UniTime Debugging}
\label{sec:simulation}

Students had to be simulated because of the following reasons:

A timetable can be created without students (only classes, rooms, instructors, preferences are needed).
You cannot \emph{evaluate} it without them, though.
The most important metric the institute cares about is whether there are any overlapping classes in a student's weekly schedule, but that isn't meaningful until students are put into sections.

Because of the lack of data about students at the institute, we used simulated students.
The simulation is not random, it is deterministic.
It begins with the actual number of students in each year, partitions students into lecture groups, assigns students to lab groups within each lecture group, and generates course requests for all required courses and one elective course per group of students.
Finally, it adds each student to a particular lecture section and lab section, creating the 9{,}148 enrolment rows in file 16.

The filling of the sections and the consequences thereof.

Enrolments are assigned one section at a time.
Each student goes into their division's preferred section until it hits the limit in the offering file, then spills into later sections of the same course and type.
This is similar to what is typically observed in over-subscribed labs, and it leaves two clear imprints on the data.

First, some lab batches are mixed, with students from two different lecture divisions being in the same batch as one division's preferred batches filled first.
This is a combination of three final year lab batches.
Second, if the section size is too small for a cohort, additional students remain unplaced and are not dropped silently from the generator, but rather logged.
That's why 200 first-year students and two M.Overflow of tech lab in Section~\ref{sec:datasets}.

These are both effects of the \emph{input data} and not the solver.
They are tracked because they are a tie that binds two divisions together in an unusual way, which can cause conflict where it is not expected.

\subsection{Conflict types encountered}

During this stage, three types of student conflict emerged and we needed to distinguish each type before we could do anything about it:

\begin{itemize}
  \item True overlap of two classes that the student is truly enrolled in. These are the types of conflicts that the solver should delete, and all 40 remaining conflicts in the final solution are of this type.
  \item Artefacts of cross-division batches, where a clash only occurs as a result of a student being placed in another division's batch by section filling.
  \item Artefacts of measurement: where the number of students reported has changed due to changes in the student population, rather than the timetable. Section~\ref{sec:experiments} analyses these.
\end{itemize}

\subsection{UniTime debugging log}

A significant amount of time was lost searching for import and solver failures that did not provide any useful error message.
This is not documented in the UniTime docs, so each symptom is listed with its root cause and fix in Table~\ref{tab:debug}.

\begin{table}[H]
\centering
\caption{Import and solver failures that were observed, with their causes and resolutions.}
\label{tab:debug}
\footnotesize
\begin{tabular}{@{}L{4.0cm}L{5.4cm}L{5.4cm}@{}}
\toprule
\textbf{Symptom} & \textbf{Cause} & \textbf{Resolution} \\
\midrule
Every class reports \emph{``has no available room (class not loaded)''} although all rooms imported
  & Buildings, built by file 7, are global and not scoped to the session, so no session has buildings
    until it is told to pick them up
  & \textbf{Administration\,$\rightarrow$\,Academic Sessions\,$\rightarrow$\,Buildings\,$\rightarrow$\,Update Data}, then reload the solver \\
Solver loads a few variables rather than hundreds
  & If there is no time pattern on the subpart, UniTime treats the class as ``arrange hours'' and excludes
    it from the solver
  & Import \texttt{13preferences.xml}, and re-import it whenever offerings are reloaded \\
The imported archive is empty
  & Data Exchange reads the archive root; \texttt{zip -r} puts the documents in a directory
  & Flatten the name of the entries, as the archive scripts do now \\
Teachers are not attached to classes
  & The staff imported by file 10 are not yet in the session's instructor list
  & \textbf{Courses\,$\rightarrow$\,Input Data\,$\rightarrow$\,Instructors\,$\rightarrow$\,Manage Instructor List} \\
Solver reports \texttt{Student conflicts: 0} and exports an empty student block
  & \emph{Projected Student Course Demands} reads last-like demand and curriculum projections;
    where the session has neither, zero students load
  & Set \emph{Student Course Demands} to \emph{Enrolled Student Course Demands} and reload \\
Re-importing offerings deletes classes unexpectedly
  & \texttt{incremental="true"} will not delete unmatched offerings, but will delete unmatched
    classes within offerings in the file
  & Uncommit any saved solution first, and treat file 12 as authoritative for the offerings it mentions \\
The exported solver XML has no class or room names
  & By default, the solution export is anonymised
  & Set \texttt{unitime.solution.export.names=true} in
    \textbf{Administration\,$\rightarrow$\,Defaults\,$\rightarrow$\,Configuration} \\
Student group import fails with a Hibernate \texttt{TransientObjectException}
  & Combining two large groups into one transaction
  & Emit group membership in the student records of file 14, rather than as bulk group updates \\
\bottomrule
\end{tabular}
\end{table}

The tendency to assign pupils to their sections.

The most significant thing we've discovered is how UniTime deals with imported enrolments.
The database loader loads each \texttt{StudentClassEnrollment} row, and attaches the student to the class provided that the student has already been attached to that course offering, and that the student is permitted to be attached to that class.
The initial sectioning still continues after the sectioning, but it does not include students who are already in the class, and the loader will only log a warning if a student left their loaded class.
To summarize: UniTime respects imported enrolments, if you tell the solver to do so.

All three are required:

\begin{enumerate}
  \item Enrolments are read at all, with \emph{Student Course Demands} set to \emph{Enrolled Student Course Demands}.
  \item \texttt{General.SwitchStudents=false}, so the last sectioning pass does not move students between parallel sections to hide conflicts.
  \item \texttt{Global.LoadStudentEnrlsFromSolution=true}, so sections are loaded from the stored enrolments and not recomputed.
\end{enumerate}

With these enabled, the conflict count rises from an apparent zero to a real 40, exactly what we wanted to see: see Section~\ref{sec:experiments}.
